\medskip

\textbf{Partie A : Détermination d'une fonction $f$ et résolution d'une équation différentielle}

\medskip



\[f(x) = \e^x+ ax + b\e^{-x}\]





\begin{center}
\psset{unit=1.25cm}
\begin{pspicture*}(-2.2,-0.6)(3.2,6.6)
\psgrid[gridlabels=0pt,subgriddiv=5,gridwidth=0.25pt,subgridwidth=0.1pt]
\psaxes[linewidth=1.25pt,labelFontSize=\scriptstyle]{->}(0,0)(-2.2,-0.4)(3.2,6.6)
\psplot[plotpoints=2000,linewidth=1.25pt,linecolor=red]{-2}{2.5}{2.71828 x exp x sub 2 2.71828 x exp div add} \uput[r](2,5){\red $\mathcal{C}$}
\psplot[linewidth=1.25pt]{-2}{2}{3 2 x mul sub}
\end{pspicture*}
\end{center}

\smallskip

\begin{enumerate}
\item ~
\starredbullet~$f(0) \approx 3$ ;

\starredbullet~$f'(0) \approx - 2$ (coefficient directeur de la tangente)
\item 
On a $f(0) = 1 + b$. Or $1 + b = 3 \iff b = 2$.
\item On admet que la fonction $f$ est dérivable sur $\R$ et on note $f'$ sa fonction dérivée. 
	\begin{enumerate}
		\item 
Sur $\R$, \, $f'(x) = \e^x + a - b\e^{-x} = \e^x + a - 2\e^{-x}$.
		\item 
Donc $f'(0) = 1 + a - 2 = a - 1$.
		\item 
		On a vu que $f'(0) = - 2 = a - 1 \iff a = - 1$, donc finalement :
		
		\[f(x) = \e^x - x + 2\e^{-x}\]
	\end{enumerate}
\item 

\[(E) : \quad  y' + y = 2\e^x - x - 1\]

	\begin{enumerate}
		\item 
		
\[g(x) = \e^x - x + 2\e^{-x}.\]

$g$ étant une somme de fonctions dérivables sur $\R$, on a sur cet intervalle :

$g'(x) = \e^x - 1 - 2\e^{-x}$, soit en remplaçant dans l'équation $(E)$ :

$\e^x - 1 - 2\e^{-x} + \e^x - x + 2\e^{-x} = 2\e^x - x - 1$ égalité vraie : $g$ est donc une solution de l'équation différentielle $(E)$.


		\item 
$y' + y = 0 \iff y' = - y$ : on sait que les solutions de cette équation différentielle sont les fonctions définies par :

\[ x \longmapsto f(x) = K \e^{- x},\, \text{avec } K \in \R.\]

		\item 
D'après les questions précédentes les solutions de l'équation $(E)$ sont les fonctions définies que $\R$ par :

\[ x \longmapsto \e^x - x + 2\e^{-x} + K \e^{- x},\, \text{avec } K \in \R.\]
	\end{enumerate}
\end{enumerate}

\bigskip

\textbf{Partie B : Étude de la fonction $g$ sur $[1~;~+\infty[$}

\medskip

\begin{enumerate}
\item 


En posant $X = \e^x$, on a $X^2 - X - 2$ et ce trinôme a deux solutions évidentes $2$ et $- 1$.

$X^2 - X - 2 = (X - 2)(X + 1)$, donc en revenant à l'écriture initiale :

$\e^{2x} - \e^x - 2 = \left(\e^x - 2\right)\left(\e^x + 1\right)$ pour tout réel $x$.
\item 
$g'(x) = \e^x - 1 - 2\e^{-x} = \e^{-x}\left(\e^{2x} - \e^{x} - 2)\right)$, donc d'après la question précédente :

$g'(x) = \e^{-x}\left(\e^x - 2\right)\left(\e^x + 1\right).$
\item 
Puisque $\e^x - 2 > 0$ (admis), $\e^x + 1 > 1 > 0$  et $\e^{-x} > 0$, donc leur produit $g'(x) > 0$.

Il en résulte que la fonction $g$ est croissante sur $\R$, donc sur $[1~;~ +\infty[$.

\end{enumerate}
